% !TeX program = xelatex
\documentclass[12pt,a4paper]{article}
\usepackage{fontspec}
\setmainfont{Times New Roman}
\usepackage[ukrainian]{babel}
\usepackage{geometry}
\geometry{left=30mm,right=15mm,top=20mm,bottom=20mm}
\usepackage{amsmath,booktabs,array,setspace,hyperref,graphicx,float}
\hypersetup{hidelinks}
\onehalfspacing
\setlength{\parindent}{1.25cm}
% Якщо зображення відсутнє, компіляція продовжиться із заглушкою.
\newcommand{\reportimage}[3]{%
	\begin{figure}[H]\centering
		\IfFileExists{#1}{\includegraphics[width=#2\textwidth,height=0.38\textheight,keepaspectratio]{#1}}{\fbox{\parbox{0.75\textwidth}{\centering Додайте файл \texttt{\detokenize{#1}} у папку зі звітом}}}
		\caption{#3}\end{figure}}
\begin{document}
	\begin{titlepage}
		\begin{center}
			\large
			МІНІСТЕРСТВО ОСВІТИ І НАУКИ УКРАЇНИ \\
			ОДЕСЬКИЙ НАЦІОНАЛЬНИЙ УНІВЕРСИТЕТ ІМЕНІ І. І. МЕЧНИКОВА \\
			Факультет математики, фізики та інформаційних технологій \\
			Кафедра оптимального керування і економічної кібернетики
			\vspace*{5cm}
		
			\Large\textbf{ЗВІТ}\\
			з лабораторної роботи №2 \\
			з дисципліни «Системи комп'ютерного зору» \\
			\vspace{0.5cm}
			\large\textbf{Тема:} «Калібрування камери та усунення дисторсії»
			\vspace{4cm}
			\begin{flushright}
				\textbf{Виконала:} \\
				студентка 3 курсу \\
				спеціальності «Прикладна математика» \\
				Овчарук Єлизавета Євгеніївна \\
				\vspace{1cm}
				\textbf{Перевірив:} \\
				доцент, к. ф.-м. н. \\
				Мороз Володимир Володимирович
			\end{flushright}
			\vfill Одеса --- 2026
		\end{center}
	\end{titlepage}
	\section{Мета роботи}
	Ознайомитися з процедурою калібрування камери за допомогою плоского шахового шаблону, визначити внутрішні параметри камери й коефіцієнти дисторсії, оцінити точність за похибкою репроєкції.
	\section{Теоретичні відомості}
	Калібрування камери --- оцінювання параметрів математичної моделі проєкції тривимірних точок на площину зображення. Матриця внутрішніх параметрів має вигляд:
	\begin{equation}
		K=\begin{bmatrix}f_x&0&c_x\\0&f_y&c_y\\0&0&1\end{bmatrix},
	\end{equation}
	де $f_x,f_y$ --- фокусні відстані у пікселях, $c_x,c_y$ --- координати головної точки. Для кожного кадру зіставляють координати вузлів шаблону в міліметрах із координатами їхніх проєкцій у пікселях. Функція \texttt{cv2.calibrateCamera} мінімізує репроєкційну похибку та оцінює параметри дисторсії.
	\section{Засоби та вихідні дані}
	\begin{itemize}
		\item Python, OpenCV, NumPy, Matplotlib.
		\item Роздільність зображень: $2560\times1920$ пікселів.
		\item Шаблон: $9\times6$ внутрішніх кутів (54 точки на кадр).
		\item Розмір сторони клітинки: 25,0 мм (за умови відповідності реальному шаблону).
		\item Основний набір: 25 знімків (10 \texttt{calib\_corner}, 5 \texttt{calib\_front}, 10 \texttt{calib\_tilt}).
	\end{itemize}
	\reportimage{calib_front_01.jpg}{0.65}{Вихідне зображення калібрувального шахового шаблону.}
	\section{Хід виконання роботи}
	\begin{enumerate}
		\item Підготовлено фотографії шахівниці з різних ракурсів.
		\item Зображення перетворено у відтінки сірого; кути знайдено методом \texttt{findChessboardCornersSB}.
		\item Сформовано відповідності двовимірних точок зображення та тривимірних точок плоского шаблону ($Z=0$).
		\item За допомогою \texttt{cv2.calibrateCamera} оцінено матрицю камери, дисторсію та положення шаблону на кадрах.
		\item Обчислено покадрову похибку репроєкції, збережено результати та побудовано графік.
	\end{enumerate}
	\reportimage{corners_calib_front_01.jpg}{0.65}{Візуалізація знайдених кутів шахівниці.}
	\section{Результати}
	За підсумковим файлом \texttt{calibration.txt} для 25 кадрів отримано:
	\begin{equation}
		K=\begin{bmatrix}
			1760.25056&0&1274.05000\\
			0&1766.42963&959.79976\\
			0&0&1
		\end{bmatrix}.
	\end{equation}
	Коефіцієнти дисторсії $[k_1,k_2,p_1,p_2,k_3]$:
	\begin{equation}
		[-0.08649133,\ 0.89210352,\ -0.00425102,\ 0.00114950,\ -1.70730945].
	\end{equation}
	Початкова та фінальна RMS: \textbf{1.436963 px}; використано \textbf{25 кадрів}.
	\reportimage{results_screen.jpg}{0.95}{Скріншот підсумкових параметрів калібрування для 25 кадрів.}
	\reportimage{reprojection_errors.png}{0.95}{Покадрова похибка репроєкції (графік для основного набору 25 кадрів).}
	\section{Аналіз отриманих результатів}
	Загальна RMS для основного набору 25 зображень становить 1.436963 px. Найбільші відхилення спостерігаються для окремих знімків серії \texttt{calib\_corner}, що може бути пов'язано зі складнішими перспективними ракурсами. Головна точка приблизно $(1274.05;959.80)$ px близька до центра кадру $(1280;960)$ px.
	
	Окремий запуск на 26 зображеннях дав RMS 1.415182 px, але включав додатковий \texttt{test\_line.jpg}; ці результати не слід змішувати з основним набором. RMS сама по собі не гарантує метричної точності. Потрібно також перевірити коректність позначення всіх 54 внутрішніх кутів та фактичний розмір клітинки.
	\section{Висновки}
	У роботі реалізовано калібрування камери засобами Python та OpenCV. Оброблено 25 фотографій шахівниці, визначено матрицю внутрішніх параметрів, коефіцієнти дисторсії та покадрові похибки. Отримана RMS 1.436963 px свідчить про працездатність алгоритму, однак остаточна оцінка якості потребує перевірки правильності відповідностей кутів та вимог до точності.
	\section{Використані джерела}
	\begin{enumerate}
		\item OpenCV Documentation. Camera Calibration: \url{https://docs.opencv.org/4.x/dc/dbb/tutorial_py_calibration.html}
		\item OpenCV Documentation. Camera Calibration and 3D Reconstruction: \url{https://docs.opencv.org/4.x/d9/d0c/group__calib3d.html}
	\end{enumerate}
	\noindent\textit{Примітка: звіт підготовлено за наданими скріншотами; вимоги викладача та оригінальні знімки не перевірялися.}
\end{document}
